\section{Question 3}
\begin{enumerate}[label=(\alph*)]
    \item 
    We can write,
    \[
    f(z) = \frac{z^2}{5(z+i)(z-i)(z^2+1/5)}
    \]
    \[
    = \frac{z^2}{5(z+i)(z-i)(z+i/\sqrt{5})(z-i/\sqrt{5})}
    \]
    and hence see that there are singularities at,
    \[
    \pm i,\ \mathrm{and}\ \pm i/\sqrt{5}
    \]
    all of which are simple poles since as $z$ approaches each of them $f\rightarrow\infty$ (by 2.\ and 3.\ HB p.\ 58 respectively).

    Now, using the Cover-Up Rule in each case (we can see that the implied function $g$ will be analytic at each singularity), we have,
    \[
    \Res(f,i) = \frac{i^2}{5(2i)(i+i/\sqrt{5})(i-i/\sqrt{5})}
    \]
    \[
    = \frac{i^2}{10i^3(1+1/\sqrt{5})(1-1/\sqrt{5})}
    \]
    \[
    = \frac{1}{10i(1-1/5)} = \frac{5}{10i(4)} = \frac{1}{8i} = -\frac{i}{8}
    \]
    and additionally,
    \[
    \Res(f,-i) = \frac{i}{8}
    \]
    since $f$ is an even function (1.\ HB p.\ 59).

    Also,
    \[
    \Res(f, i/\sqrt{5}) = \frac{i^2/5}{5(i/\sqrt{5} - i)(i/\sqrt{5} + i)(2i/\sqrt{5})}
    \]
    \[
    = \frac{i^2/5}{2\sqrt{5}i^3(1/\sqrt{5} - 1)(1/\sqrt{5} + 1)}
    \]
    \[
    = \frac{1/5}{2\sqrt{5}i(1/5 - 1)}
    \]
    \[
    = \frac{1}{2\sqrt{5}i(1-5)} = -\frac{1}{8\sqrt{5}i} = \frac{i}{8\sqrt{5}}
    \]
    and therefore,
    \[
    \Res(f,-i/\sqrt{5}) = -\frac{i}{8\sqrt{5}}
    \]
    for the same reason.

    \item 
    By 8.\ HB p.\ 62, we have,
    \[
    \frac{p(t)}{q(t)} = \frac{t^2}{(t^2+1)(5t^2+1)}
    \]
    where the degree of $q$ is $4$ which exceeds that of $p$ by $2$ (the degree of $p$ is $2$), and there are no poles on the real axis.

    Then,
    \[
    \int_{-\infty}^{\infty} \frac{p(t)}{q(t)}\dd t = 2\pi i S + \pi i T
    \]
    \[
    = 2\pi iS
    \]
    where $S$ is the sum of the residues of $p/q$ at the poles in the upper half plane, and $T$ is the sum of poles on the real axis.

    We have from (a),
    \[
    S = \Res(f,i) + \Res(f,i/\sqrt{5}) = -\frac{i}{8} + \frac{i}{8\sqrt{5}}
    \]
    \[
    = \frac{i}{8}\left(\frac{1}{\sqrt{5}} - 1\right) = \frac{i}{8}\left(\frac{1-\sqrt{5}}{\sqrt{5}}\right) = \frac{i(1-\sqrt{5})\sqrt{5}}{40} = \frac{i(\sqrt{5}-5)}{40}
    \]
    Therefore the solution to the integral is,
    \[
    \int_{-\infty}^{\infty} \frac{p(t)}{q(t)}\dd t = -\frac{2\pi(\sqrt{5}-5)}{40} = \frac{\pi(5 - \sqrt{5})}{20}
    \]

    \textbf{TN: .... Confirm this is correct!}

    {
    \color{blue}

    I agree with your answer for the whole of this. I arrived at the same residues using the $g/h$ rule, and then used the rule for even functions too.

    Although, in the heat of the moment, I made a stupid slip at the very end for the contour integral. I remember spending so long getting the factorisation of the numerator right that I didn't cancel properly, and $2\pi/40 \cdots$ became $\pi/40 \cdots$!

    John
    }
    
    \textcolor{purple}{Thanks for confirming that John. Ah, I have noticed a lot of my own mistakes by putting together these questions! On the day, I had an imaginary answer for the final real integral... But I was so limited in time I didn't go back to correct it!
    Well, we move on with life, maths isn't just about exams after all.}
    
\end{enumerate}
